% !TeX root = ../main.tex
%
% Compute statement.
%
% ANONYMITY: naming the cloud provider and the VM SKU is fine.  Naming the
% workspace host, the account path or the repository is not.  Do not add them.
%
% The archived-record count below (97) is the number of run records in the
% pre-protocol archive snapshot, counted from the archive itself.  If more
% pre-protocol runs are archived before submission, update it -- do not guess.

\section{Compute}
\label{sec:stmt:compute}

\paragraph{Hardware.}
Every run reported in this paper was executed on a single NVIDIA A100 GPU, on a
cloud node with 24 vCPUs (an Azure \texttt{Standard\_NC24ads\_A100\_v4}
instance). Runs are serialised: one training job owns the GPU at a time and is
given all 24 data-loading workers. There is no multi-GPU or multi-node training
anywhere in this work, and no result depends on a distributed configuration.

\paragraph{The reported sweep.}
The matched grid is four backbones $\times$ three pruning criteria $\times$ four
sparsities $\times$ two arms, plus one dense baseline per backbone: 100 cells.
Executed end to end on the machine above, at one seed per cell, the sweep takes
roughly 14 hours of wall-clock. Replicating the headline cells at three seeds
costs two further passes over that subset. Individual cells are not uniform in
cost: a dense pre-training run has the largest epoch budget, and a contrastive
run is materially more expensive per epoch than the baseline it is compared
against, because it forwards two augmented views through the student and a
frozen teacher rather than one view through a single network. The matched-budget
claim in this paper is a claim about \emph{epochs and schedule}, which are
identical across arms; it is not a claim that the two arms cost the same number
of GPU-seconds, and they do not.

\paragraph{Supporting runs reported in the paper.}
Beyond the grid, the paper reports three smaller sets of runs. The objective
ablation is nine configurations --- three sparsities crossed with three
objective variants --- of which eight completed; one was terminated before it
finished and is reported as not run rather than as a gap. The learning-rate
measurement for the VGG family is three runs at one model, criterion and
sparsity. The run-to-run noise floor was estimated from a five-seed repetition
of one dense configuration.

\paragraph{Compute this project consumed that this paper does not report.}
Considerably more. A body of 97 earlier run records exists and is archived
rather than used. Those runs are excluded for a protocol reason, not for their
results: they were produced under a superseded configuration --- a different
epoch budget, a different mask-update interval, and a contrastive stage that
trained on the unsplit 50{,}000-image training set rather than on the 45{,}000
that the current protocol reserves for it --- so they are not comparable with the
matched sweep and mixing them into it would silently contaminate a matched
table with an unmatched number. They are archived intact rather than deleted.
Separately, an earlier line of work on this codebase produced results that were
invalidated by an evaluation defect, which is described in the reproducibility
statement; those numbers are not reported anywhere in this paper. Smoke and
tier-zero runs, used to exercise the pipeline on truncated data, are excluded
mechanically by the results generator and are not counted as experiments. We do
not report a total GPU-hour figure for the discarded work, because the archived
records span several machines and configurations and any single figure we gave
would be an estimate rather than a measurement.

\paragraph{Unmeasured cells.}
Cells that had not been measured at the time of writing appear in this document
as a visible placeholder rather than as a number, and are not counted anywhere
as results.
